\documentclass[UTF8,a4paper,zihao=-4]{ctexart}

\usepackage[a4paper,margin=2.4cm]{geometry}
\usepackage{booktabs}
\usepackage{amsmath}
\usepackage{array}
\usepackage{algorithm}
\usepackage{algpseudocode}
\usepackage[numbers,sort&compress]{natbib}
\usepackage[hidelinks]{hyperref}

\floatname{algorithm}{算法}
\renewcommand{\abstractname}{摘要}
\renewcommand{\refname}{参考文献}
\algrenewcommand\algorithmicrequire{\textbf{输入：}}
\algrenewcommand\algorithmicensure{\textbf{输出：}}

\title{面向工人柔性的柔性作业车间调度：\\工人感知关键路径禁忌搜索}
\author{匿名作者}
\date{\today}

\begin{document}
\maketitle

\begin{abstract}
经典柔性作业车间调度问题（FJSSP）主要描述机器柔性，却通常忽略工人资源。
在真实人机协同车间中，一道工序只有同时选择兼容机器和兼容工人才可加工，并且加工时间由工序、机器、工人三元组共同决定。
带工人柔性的 FJSSP（FJSSP-W）或双资源约束 FJSP（DRCFJSP）因此引入第三类决策变量和第三类资源冲突边。
这意味着局部搜索不能只在机器析取图上移动工序，再事后修复工人冲突；析取图、邻域动作、近似评价和合法性判断都必须显式感知工人资源。
本文给出一套面向比赛场景的工人感知关键路径禁忌搜索方法。
该方法继承强单解 FJSP 局部搜索的轨迹搜索思想，但将表示、析取图、邻域、合法性检查和提交级结果生成流程全面扩展到机器与工人双资源。
在当前 30 个 Scenario 1 官方算例上，最新审计通过的结果相对于参考提交在 11 个算例上改进最优 makespan，在其余 19 个算例上保持最优值不退化；30 个算例的最优值总和从 38496 降至 38425，十次运行平均值总和从 39549.3 降至 38464.3。
\end{abstract}

\noindent\textbf{关键词：} 柔性作业车间调度；工人柔性；双资源调度；禁忌搜索；局部搜索；热启动

\section{引言}
柔性作业车间调度问题（FJSSP）是制造系统调度中的经典 NP-hard 问题。
在该问题中，每道工序可以从若干候选机器中选择一台进行加工，目标通常是最小化最大完工时间。
这一抽象对机器柔性建模很有效，但对真实人机协同车间仍然不够。
在实际生产中，机器往往需要由具备相应技能的工人操作；同一台机器在不同工人操作下也可能具有不同加工时间。
如果忽略工人资源，得到的排程可能只是在机器维度上可行，在真实车间中仍然会因为工人冲突而不可执行。

FJSSP-W 进一步引入工人柔性：每道工序不仅需要选择机器，还需要选择兼容工人；加工时间不再只依赖机器，而是依赖工序、机器和工人的组合~\cite{gong2018workerflex,gong2020habc,li2025ceamcp,hutter2025benchmark,hutter2024cec}。
在 DRCFJSP 文献中，这一问题也被描述为机器分配、工人分配和工序排序三个子问题的联合优化。
这个扩展看起来只是多了一类资源，但对局部搜索的影响很大。
标准 FJSSP 的主要冲突来自工件顺序和机器占用，而 FJSSP-W 同时存在工人占用冲突。
如果搜索仍然只看机器边，那么一个表面上缩短加工时间的移动，可能会把工序插入拥堵的工人队列，最终使真实 makespan 变差，甚至产生不可行排程。

近年不少研究进一步加入绿色指标、工人学习、疲劳、分布式工厂等扩展约束，使问题越来越贴近复杂制造场景。
相比之下，FJSSP-W 比赛保留的是最核心的双资源可行性和 makespan 目标。
这使它很适合验证一个基础问题：标准 FJSP 中强大的关键路径局部搜索，是否可以升级为真正的双资源局部搜索。

本文的算法受 Zhang 等人强单解 FJSP 局部搜索思想启发~\cite{zhang2024awls}，但本文叙事重点不是“修改某个已有算法”，而是给出 worker-aware 的升级方式。
我们的目标不是重新设计一个种群式元启发式，而是把机器感知的局部搜索变成同时感知机器和工人的轨迹搜索。

本文贡献可以概括为四点：
\begin{itemize}
  \item 给出面向 FJSSP-W 的三元解表示，将工序序列、机器分配和工人分配同时纳入搜索状态；
  \item 在析取图中加入工人边，使关键结构、合法性判断和近似评价都能看到工人资源；
  \item 设计包含顺序移动、机器重分配、工人重分配和机器工人联合重分配的混合邻域；
  \item 建立热启动探索、十种子提交级再生成和审计晋级流程，避免把单个幸运结果误当成合规提交。
\end{itemize}

\section{问题定义}
设 $O$ 为全部工序集合。
每道工序 $i\in O$ 属于某个工件，因此同一工件内部工序之间存在固定先后关系。
对每道工序，允许的加工方案是若干个二元组 $(m,w)$，其中 $m$ 是可行机器，$w$ 是可行工人。
加工时间记为 $p_{i,m,w}$；如果三元组不可行，则对应加工时间不存在。

一个完整排程需要为每道工序确定开始时间 $s_i$、机器 $m_i$ 和工人 $w_i$。
排程必须满足以下约束：
\begin{enumerate}
  \item 工件内部的工序先后约束；
  \item 同一机器上的工序不能时间重叠；
  \item 同一工人处理的工序不能时间重叠；
  \item 每道工序所选机器工人组合必须可行，即 $p_{i,m_i,w_i}>0$。
\end{enumerate}
优化目标是最小化最大完工时间：
\[
  C_{\max}=\max_{i\in O}\left(s_i+p_{i,m_i,w_i}\right).
\]

\section{算法}
\subsection{从机器单资源搜索到双资源搜索}
经典 FJSSP 局部搜索面向标准 FJSSP。
在该问题中，一个移动可以抽象为“把工序 $i$ 放到机器 $m$ 的位置 $k$”。
在 AWLS 这一类强单解局部搜索中，算法维护单个当前解，围绕关键路径生成候选移动，用禁忌表避免短期循环，并用快速 makespan 估计来筛选邻域。
这一类方法最有辨识度的机制之一是自适应权重：每个工序维护累计权重和空闲计数，近似评价时把权重转换为额外加工时间。
当某次移动未能改进当前解时，被移动工序在之后的估计中会变得“不那么有吸引力”；当出现新的全局最好解时，权重被重置。
因此，禁忌表提供短期记忆，自适应权重提供长期扰动，两者共同帮助搜索跳出局部最优。

FJSSP-W 使得原版移动定义不再充分。
一个移动不仅要说明工序放在哪台机器上，还要说明由哪个工人加工，以及它在对应工人队列中的相对位置。
因此，我们保留单解禁忌搜索框架，但把搜索状态、析取图、邻域、近似评分和精确验证全部扩展到工人维度。
当前提交级结果使用的是保守的工人感知版本：以工人边和精确解码为合法性底线，而不是把完整自适应权重直接作为当前结果的来源。
完整 AWLS 权重项仍是下一步值得重新开启和调参的消融方向。

\subsection{解表示与解码}
算法用三部分表示一个解：
\[
  x=(\pi,M,W),
\]
其中 $\pi$ 是工件序列表示，$M$ 给出每道工序的机器选择，$W$ 给出每道工序的工人选择。
如果工件 $j$ 在 $\pi$ 中第 $k$ 次出现，则该位置代表工件 $j$ 的第 $k$ 道工序。
这样，工件内部先后关系由出现次数隐式保证，而 $M_i$ 和 $W_i$ 决定工序 $i$ 的具体加工组合。

解码器从左到右扫描 $\pi$。
当遇到工件 $j$ 的下一道工序 $i$ 时，开始时间取三个 ready time 的最大值：
\[
  s_i=\max\{R^{job}_j,\ R^{mach}_{M_i},\ R^{work}_{W_i}\}.
\]
其中 $R^{job}_j$ 是工件 $j$ 上一道工序的完工时间，$R^{mach}_{M_i}$ 是所选机器当前可用时间，$R^{work}_{W_i}$ 是所选工人当前可用时间。
调度完工序 $i$ 后，这三个 ready time 都用 $p_{i,M_i,W_i}$ 更新。
只要每个 $(M_i,W_i)$ 都在可行集合中，这个左移解码器就能产生一个可行主动排程。

\subsection{加入工人边的析取图}
标准 FJSSP 的局部搜索通常基于析取图理解排程。
图中包含工件先后边和机器析取边。
在 FJSSP-W 中，我们加入第三类边：工人边。
对一个已解码排程，构造如下边集合：
\begin{itemize}
  \item 工件边：同一工件相邻工序之间的先后关系；
  \item 机器边：同一机器队列中相邻工序之间的先后关系；
  \item 工人边：同一工人队列中相邻工序之间的先后关系。
\end{itemize}

工人边不是装饰性的。
它决定了工人重分配之后是否可行，也影响关键路径判断。
在这个图中，关键工序是至少位于一条源点到汇点最长路径上的工序。
由于最长路径现在可能通过机器边，也可能通过工人边，所以即使机器队列完全不变，一个纯工人移动也可能改变关键结构。

\subsection{精确合法性验证}
所有被接受的移动都必须经过精确验证。
验证包括三步：
\begin{enumerate}
  \item 检查被修改工序的机器工人组合是否在预计算可行列表中；
  \item 用同时包含工件、机器和工人 ready time 的解码器重新计算排程；
  \item 根据解码结果重建机器队列和工人队列，检查是否存在重叠。
\end{enumerate}
第三步看似冗余，但对比赛代码非常重要。
它可以捕捉近似评价或增量更新中的实现错误，防止不可行排程进入最终 CSV。

\subsection{邻域设计}
搜索使用混合邻域。
顺序邻域包括交换、插入、相邻交换、块交换、块迁移、双移动和反转。
分配邻域包括机器重分配、工人重分配、机器相邻交换、工人相邻交换，以及机器工人联合重分配。
所有分配移动都先经过工序级可行列表过滤，不可行的机器工人组合不会进入精确评价。

邻域采样偏向当前关键结构。
对顺序移动，算法优先选择位于长机器队列或长工人队列附近的工序，因为这些队列更可能决定 $C_{\max}$。
对分配移动，算法从该工序的可行组合中选择候选，并优先考虑加工时间较短或目标队列局部压力较小的组合。
与标准 FJSSP 最大的区别是：重分配不再只是机器决策。
如果只改变 $M_i$ 而不重新考虑 $W_i$，瓶颈可能只是从机器队列转移到工人队列。

从概念上说，原版 AWLS 的移动 $(i,m,k)$ 被扩展为两类移动。
资源移动改变 $(M_i,W_i)$；顺序移动在机器队列、工人队列或二者共同诱导的局部结构中改变工序顺序。
对联合分配移动，算法先选择可行组合 $(m,w)\in E_i$，再测试由当前机器队列和工人队列决定的插入位置。
这样可以减少大量明显不可能的三元组合，使有限的精确评价预算集中在更可能可行的候选上。

\subsection{禁忌搜索与近似筛选}
主搜索是带 aspiration 的禁忌搜索。
每类移动都有相应的禁忌键，用来记录被改变的顺序或分配属性。
每次迭代中，算法采样一批候选移动，先用快速筛选层评分，再只对排名靠前的 $K$ 个候选做精确解码。
近似筛选只负责排序，不负责最终接受。

近似评分继承了 AWLS 中“用估计值筛选邻域”的思想，但从单一机器队列扩展为机器和工人两个资源队列。
对影响工序 $i$ 的候选移动 $a$，评分可以写成：
\[
  \hat{C}(a)=\Delta p_i(a)+\rho_m(a)+\rho_w(a)+\lambda Z_i .
\]
其中 $\Delta p_i(a)$ 是加工时间变化，$\rho_m(a)$ 是目标机器局部队列压力，$\rho_w(a)$ 是目标工人局部队列压力，$Z_i$ 是由 AWLS 工序权重和空闲计数诱导的额外加工时间。
当前保守版本取 $\lambda=0$，即不把完整 AWLS 权重项计入当前提交级结果。
保留该公式的原因是说明：原版 AWLS 自适应权重可以自然接入工人感知估计器，但需要单独调参和消融验证。

算法~\ref{alg:search-zh} 给出整体流程。
禁忌表记录已接受移动的反向属性，例如工序原位置或原机器工人组合。
如果某个移动处于禁忌状态，只有当它满足 aspiration 条件，即能够改进当前运行的历史最优值时，才允许被接受。
如果没有非禁忌改进移动，算法也会接受损失最小的合法候选，以维持轨迹移动并穿过平台区。

\begin{algorithm}[t]
  \caption{工人感知热启动禁忌搜索}
  \label{alg:search-zh}
  \begin{algorithmic}[1]
    \Require 算例 $I$，随机种子 $r$，初始解或热启动解 $x_0$，预算 $B$
    \Ensure 最优可行解 $x^\star$
    \State 为每道工序 $i$ 构造可行列表 $E_i=\{(m,w,p_{i,m,w})\}$
    \State $x\gets$ 修复并解码 $(x_0,E)$；若失败则随机构造可行解
    \State $x^\star\gets x$；初始化禁忌表 $T\gets\emptyset$ 和可选 AWLS 权重 $Z_i\gets0$
    \While{预算 $B$ 未耗尽}
      \State 围绕 $x$ 采样顺序、机器、工人和联合分配移动
      \ForAll{候选移动 $a$}
        \If{$a$ 选择了不可行机器工人组合}
          \State 丢弃 $a$
        \ElsIf{$a$ 为禁忌且不满足 aspiration}
          \State 丢弃 $a$
        \Else
          \State 用加工时间变化、机器压力、工人压力和可选 $Z_i$ 对 $a$ 评分
        \EndIf
      \EndFor
      \State 选取评分最好的 $K$ 个移动进入精确验证
      \State 对每个候选用工件、机器、工人 ready time 精确解码
      \State 按接受准则选择最好的合法候选 $y$
      \If{没有候选通过精确验证}
        \State 从 $x^\star$ 重启或重新采样
      \Else
        \State $x\gets y$；将反向移动属性加入 $T$；更新可选 AWLS 权重
        \If{$C_{\max}(x)<C_{\max}(x^\star)$}
          \State $x^\star\gets x$
        \EndIf
      \EndIf
    \EndWhile
    \State \Return $x^\star$
  \end{algorithmic}
\end{algorithm}

\subsection{为什么近似评价不会破坏可行性}
近似评价只用于决定“哪些候选值得精确解码”，不会直接把解写入 incumbent 或提交文件。
这一点在 FJSSP-W 中比在标准 FJSSP 中更关键。
一个移动可能缩短某道工序的加工时间，却把它插入拥堵的工人队列，从而使真实 makespan 变差。
因此，安全边界是：近似评价负责筛选候选，精确解码负责决定候选是否真实可行以及是否值得接受。

\subsection{热启动与提交晋级流程}
比赛流程本身也是算法的一部分。
我们把结果分成三个层级：
\begin{enumerate}
  \item \textbf{探索扫描：} 用较短时间和少量种子在目标算例上搜索，允许内部发现单行更优解；
  \item \textbf{提交级再生成：} 如果探索扫描发现严格改进，则用该解作为热启动，对该算例重新生成十个不同种子结果；
  \item \textbf{晋级审计：} 只有当完整 30 算例 CSV 每个官方算例恰好包含 10 行，并且相对参考文件没有最优值退化时，才允许晋级为候选提交。
\end{enumerate}
这个流程避免了一个常见陷阱：把某个幸运单次结果直接混入 CSV，然后误以为它满足了多次运行提交规范。

\section{实验设置}
\subsection{算例与参考结果}
实验使用 FJSSP-W benchmarking 环境和 Scenario 1 提交格式。
候选文件与参考提交 \texttt{atpts\_final\_submission.csv} 比较。
当前审计通过的候选为 2026-04-27 的 \texttt{best30\_v18} 结果。
两个文件都包含 30 个官方算例，每个算例恰好 10 行。

\subsection{运行规范}
提交级再生成使用十个不同随机种子：
\[
7,19,29,41,53,67,79,97,131,173.
\]
探索扫描使用较小种子批次和实例级时间限制。
一旦发现改进，再用提交级十种子重新生成目标算例。
主要第二阶段晋级结果使用 20000 次精确评价上限、200000 次迭代上限和邻域规模 60。

\subsection{评价指标}
我们报告每个算例的最优 makespan、十行平均 makespan，以及可用时达到最优行的函数评价次数。
比较类别如下：
\begin{itemize}
  \item \textbf{严格改进：} 候选最优 makespan 小于参考最优 makespan；
  \item \textbf{同最优但平均值更好：} 最优值相同，但十次平均 makespan 更小；
  \item \textbf{同最优且函数评价更少：} 最优值和平均值相同，但平均函数评价次数更低。
\end{itemize}

\section{结果}
表~\ref{tab:summary-zh} 给出 30 个 Scenario 1 算例的审计汇总。
候选结果在 11 个算例上改进参考最优 makespan，在其余 19 个算例上保持最优值相同，没有任何最优值退化。

\begin{table}[t]
  \centering
  \caption{30 个 Scenario 1 算例的审计汇总。}
  \label{tab:summary-zh}
  \begin{tabular}{lr}
    \toprule
    类别 & 数量 \\
    \midrule
    严格最优值改进 & 11 \\
    同最优但平均 makespan 更低 & 14 \\
    同最优但函数评价次数更低 & 5 \\
    最优值退化 & 0 \\
    \midrule
    官方算例数 & 30 \\
    每个算例行数 & 10 \\
    \bottomrule
  \end{tabular}
\end{table}

表~\ref{tab:strict-zh} 列出所有严格改进算例。
最大的单算例改进主要出现在 DPpaulli 和 Hurink vdata 系列。
30 个算例参考最优值总和为 38496，候选结果降至 38425，总下降 71。
十次运行平均 makespan 总和从 39549.3 降至 38464.3。

\begin{table*}[t]
  \centering
  \caption{候选结果相对参考最优 makespan 的严格改进。}
  \label{tab:strict-zh}
  \begin{tabular}{lrrrrr}
    \toprule
    算例 & 参考最优 & 候选最优 & 差值 & 参考平均 & 候选平均 \\
    \midrule
    \texttt{3\_DPpaulli\_15\_workers} & 2213 & 2199 & -14 & 2241.1 & 2208.9 \\
    \texttt{2d\_Hurink\_vdata\_30\_workers} & 1075 & 1064 & -11 & 1084.3 & 1068.1 \\
    \texttt{3\_DPpaulli\_18\_workers} & 2215 & 2204 & -11 & 2242.0 & 2210.6 \\
    \texttt{3\_DPpaulli\_9\_workers} & 2087 & 2077 & -10 & 2099.2 & 2083.5 \\
    \texttt{2c\_Hurink\_rdata\_38\_workers} & 1539 & 1533 & -6 & 1556.8 & 1536.6 \\
    \texttt{2c\_Hurink\_rdata\_28\_workers} & 789 & 784 & -5 & 811.9 & 787.0 \\
    \texttt{2d\_Hurink\_vdata\_18\_workers} & 1041 & 1036 & -5 & 1047.5 & 1036.0 \\
    \texttt{6\_Fattahi\_14\_workers} & 546 & 543 & -3 & 570.6 & 544.0 \\
    \texttt{0\_BehnkeGeiger\_60\_workers} & 403 & 401 & -2 & 419.6 & 402.0 \\
    \texttt{3\_DPpaulli\_1\_workers} & 2482 & 2480 & -2 & 2528.0 & 2481.8 \\
    \texttt{6\_Fattahi\_20\_workers} & 1155 & 1153 & -2 & 1177.5 & 1154.8 \\
    \bottomrule
  \end{tabular}
\end{table*}

\subsection{第二阶段发现}
两个最近的改进说明了有界探索流程的价值。
首先，对 Hurink vdata 30 工人算例的有界扫描发现了 1069 的热启动解；提交级十种子再生成后进一步得到 1064。
其次，对 Hurink vdata 18 工人算例的扫描发现了 1036；再生成后十行全部达到 1036，并且从该热启动出发不需要额外函数评价即可保留该值。
两个改进都通过了完整审计。

\section{讨论}
这些结果不是简单替换幸运单行得到的。
每个晋级 CSV 都保持完整 30 算例形状，每个官方算例恰好 10 行，并且相对参考提交检查最优值退化。
这很重要，因为很多最强改进先在探索模式中出现，随后才被转换成合规十种子结果。

方法上的核心经验是：FJSSP-W 必须显式看见工人资源。
在许多算例中，工人队列和机器队列一样可能成为瓶颈。
如果邻域只检查机器冲突，近似评价可能接受一些表面有利但最终被工人可行性抵消的移动。
加入工人边和工人邻接信息后，搜索拥有更接近真实排程结构的局部模型。

本文仍有两个限制。
第一，目前报告的是比赛工程流程，而不是完整消融实验。
各类邻域、工人感知筛选项和原版 AWLS 自适应权重项的边际贡献还需要单独分离。
第二，参考比较基于十行提交文件，而不是每个算例 20 到 30 次独立重复的统计实验。
如果要形成正式投稿版本，这部分还需要补齐。

\section{结论}
本文给出了一套面向 FJSSP-W 的工人感知热启动禁忌搜索方法。
该方法继承 AWLS 的单解局部搜索思想，并将搜索状态、析取图、邻域和合法性验证扩展到机器与工人双资源约束。
在当前 30 个 Scenario 1 算例上，该方法改进 11 个参考最优值，另外 19 个保持最优值相同，并在没有最优值退化的前提下降低了最优值总和和平均值总和。
后续工作应重新开启并调优完整 AWLS 自适应工序权重项，补充统计实验，并对工人边、分配移动、近似筛选和热启动晋级流程进行消融分析。

\bibliographystyle{ACM-Reference-Format}
\bibliography{references}

\end{document}
